\documentclass[11pt,a4paper]{article}
\usepackage{amsmath}
\usepackage{amsfonts}
\usepackage{amssymb}
\usepackage{graphicx}
\usepackage{float}

\title{A Method for Computing Inverse Parametric PDE Problems with Random-Weight Neural Networks}
\author{Summary of Dong and Wang (2023)}
\date{}

\begin{document}

\maketitle

\section{Introduction}

This paper presents a method for solving inverse parametric partial differential equations (PDEs) using randomized neural networks, extending the local extreme learning machine (locELM) technique to inverse problems. The method focuses on simultaneously determining unknown parameters (as constants or field distributions) and the solution field to parametric PDEs based on sparse and potentially noisy measurement data.

\section{Algorithm Description}

\subsection{Problem Formulation}

The method addresses inverse problems described by parametric PDEs of the form:
\begin{align}
\alpha_1 L_1(u) + \alpha_2 L_2(u) + \cdots + \alpha_n L_n(u) + F(u) &= f(x), \quad x \in \Omega \\
Bu(x) &= g(x), \quad x \in \partial\Omega \\
Mu(\xi) &= S(\xi), \quad \xi \in \Omega_s \subset \Omega
\end{align}

where:
\begin{itemize}
    \item $L_i$ and $F$ are differential or algebraic operators (linear or nonlinear)
    \item $u(x)$ is the unknown scalar field to be determined
    \item $\alpha_i$ are $n$ unknown parameters (constants or field distributions)
    \item $B$ is a linear operator representing boundary conditions
    \item $M$ is an operator representing measurement operations
    \item $S(\xi)$ is the measurement data at points $\xi \in \Omega_s$
\end{itemize}

\subsection{Local Randomized Neural Networks}

The domain $\Omega$ is decomposed into non-overlapping sub-domains $\Omega = \Omega_1 \cup \Omega_2 \cup \cdots \cup \Omega_N$, and the solution field $u(x)$ is represented by:

\begin{equation}
u(x) = \begin{cases}
u_1(x), & x \in \Omega_1 \\
u_2(x), & x \in \Omega_2 \\
\vdots \\
u_N(x), & x \in \Omega_N
\end{cases}
\end{equation}

On the interior sub-domain boundaries, $C^k$ continuity conditions are imposed on $u(x)$, where $k$ relates to the order of the PDE.

On each sub-domain $\Omega_i$, a local feedforward neural network (FNN) is employed with the following characteristics:
\begin{itemize}
    \item Input layer: $d$ nodes representing the input coordinates $x = (x_1, x_2, \ldots, x_d)$
    \item Output layer: Single node representing $u_i(x)$
    \item Hidden layers: $(L-1)$ hidden layers in between
    \item Architectural vector: $M = [m_0, m_1, \ldots, m_L]$ where $m_0 = d$ and $m_L = 1$
    \item The weights/biases in all hidden layers are randomly assigned on $[-R_m, R_m]$ and fixed (not trainable)
    \item Only the output-layer weights are trainable parameters
\end{itemize}

For each sub-domain $\Omega_i$, the solution is represented as:
\begin{equation}
u_i(x) = \sum_{j=1}^{M} \beta_{ij} \phi_{ij}(x) = \Phi_i(x) \beta_i
\end{equation}

where $M = m_{L-1}$ is the width of the last hidden layer, $\phi_{ij}(x)$ are the output fields of the last hidden layer, and $\beta_{ij}$ are the trainable output-layer coefficients.

\subsection{Three Training Algorithms}

The paper develops three algorithms for training the neural network:

\subsubsection{NLLSQ Algorithm}

The Nonlinear Least Squares (NLLSQ) algorithm determines the inverse parameters $\alpha$ and the trainable parameters $\beta$ simultaneously using a nonlinear least squares method with perturbations (NLLSQ-perturb).

The residual vector is defined as:
\begin{equation}
R(\alpha, \beta) = \begin{bmatrix}
R_{pde}(\alpha, \beta) \\
R_{bc}(\beta) \\
R_{mea}(\beta) \\
R_{ck}(\beta)
\end{bmatrix}
\end{equation}

This system is solved using nonlinear least squares combined with a perturbation strategy to avoid being trapped in local minima. The key features of this algorithm include:
\begin{itemize}
    \item Initial guess perturbation: If the returned solution exceeds a tolerance, a sub-iteration procedure is triggered with new initial guesses
    \item Forward-mode auto-differentiation is used for computing derivatives, which is crucial for performance
    \item Measurement residual scaling: For noisy data, scaling the measurement residual vector by a factor $\lambda_{mea}$ can improve accuracy
\end{itemize}

\subsubsection{VarPro-F1 Algorithm}

The Variable Projection Formulation 1 (VarPro-F1) algorithm eliminates the inverse parameters $\alpha$ from the problem to attain a reduced problem about $\beta$ only:

\begin{itemize}
    \item The system is rearranged into $H(\beta)\alpha = b(\beta)$
    \item For any given $\beta$, the least squares solution for $\alpha$ is $\alpha = H^+(\beta)b(\beta)$
    \item Substituting this expression gives a reduced residual $r(\beta) = H(\beta)H^+(\beta)b(\beta) - b(\beta)$
    \item The optimal $\beta^*$ is determined by minimizing $\frac{1}{2}||r(\beta)||^2$
    \item Once $\beta$ is determined, $\alpha$ is computed by linear least squares
\end{itemize}

\subsubsection{VarPro-F2 Algorithm}

The Variable Projection Formulation 2 (VarPro-F2) algorithm eliminates the trainable network parameters $\beta$ from the problem to attain a reduced problem about $\alpha$ only:

\begin{itemize}
    \item The system is rearranged into $H(\alpha)\beta = b$
    \item For any given $\alpha$, the least squares solution for $\beta$ is $\beta = H^+(\alpha)b$
    \item Substituting this expression gives a reduced residual $r(\alpha) = H(\alpha)H^+(\alpha)b - b$
    \item The optimal $\alpha^*$ is determined by minimizing $\frac{1}{2}||r(\alpha)||^2$
    \item Once $\alpha$ is determined, $\beta$ is computed by linear least squares
\end{itemize}

For PDEs that are nonlinear with respect to the field solution, VarPro-F2 needs to be combined with a Newton iteration.

\section{Theoretical Findings}

\subsection{Exponential Convergence Properties}

A significant theoretical finding is that for smooth solutions and noise-free data, the errors for the inverse parameters and the solution field decrease exponentially with respect to:
\begin{itemize}
    \item The number of collocation points
    \item The number of training parameters
\end{itemize}

This exponential convergence is analytically demonstrated and empirically verified across multiple test cases. For sufficiently large parameter values, the errors can reach a level close to machine accuracy.

\subsection{Random-Weight Neural Networks as Universal Approximators}

The method leverages the proven theoretical results that single-hidden-layer feedforward neural networks with random but fixed hidden units can approximate any continuous function to any desired degree of accuracy, provided that the number of hidden units is sufficiently large.

\subsection{Properties of the Variable Projection Formulation}

The theoretical benefits of the variable projection approach include:
\begin{itemize}
    \item Reduced dimension of parameter space
    \item Better conditioning of the optimization problem
    \item Faster convergence with the reduced problem
\end{itemize}

The VarPro-F1 and VarPro-F2 algorithms are reciprocal formulations of the variable projection idea, with the former being more suitable for general cases and the latter being more efficient for linear (with respect to solution field) parametric PDEs.

\section{Numerical Examples and Results}

The method was rigorously tested using several inverse parametric PDE problems:

\subsection{Parametric Poisson Equation}

The inverse problem of determining the coefficient $\alpha$ in:
\begin{align}
\frac{\partial^2 u}{\partial x^2} + \alpha \frac{\partial^2 u}{\partial y^2} &= f(x,y) \\
\text{with boundary conditions and measurement data}
\end{align}

Key findings:
\begin{itemize}
    \item For noise-free data, the maximum error for $u$ was on the order of $10^{-7}$ in the domain
    \item The relative error of the computed $\alpha$ was approximately $9.03 \times 10^{-9}$
    \item Exponential convergence with respect to the number of collocation points
    \item Exponential convergence with respect to the number of training parameters
    \item With 10\% noise in the measurement data, the relative error of $\alpha$ was around 1\%
\end{itemize}

\subsection{Parametric Nonlinear Helmholtz Equation}

The inverse problem of determining coefficients $\alpha_1$ and $\alpha_2$ in:
\begin{align}
\frac{\partial^2 u}{\partial x^2} + \frac{\partial^2 u}{\partial y^2} - \alpha_1 u + \alpha_2 \cos(2u) &= f(x,y) \\
\text{with boundary conditions and measurement data}
\end{align}

Key findings:
\begin{itemize}
    \item VarPro-F1 exhibited high accuracy with maximum error on the order of $10^{-8}$
    \item The relative errors of $\alpha_1$ and $\alpha_2$ were $3.52 \times 10^{-11}$ and $4.76 \times 10^{-10}$
    \item Exponential convergence with respect to collocation points and training parameters
    \item With 1\% noise, relative errors for $\alpha_1$ and $\alpha_2$ were on the order of 0.1\%
\end{itemize}

\subsection{Parametric Viscous Burgers' Equation}

The inverse problem of determining $\alpha_1$ and $\alpha_2$ in:
\begin{align}
\frac{\partial u}{\partial t} + \alpha_1 u \frac{\partial u}{\partial x} &= \alpha_2 \frac{\partial^2 u}{\partial x^2} + f(x,t) \\
\text{with boundary/initial conditions and measurement data}
\end{align}

Key findings:
\begin{itemize}
    \item The maximum error for $u$ was on the order of $10^{-8}$ in the domain
    \item The relative errors of $\alpha_1$ and $\alpha_2$ were $1.30 \times 10^{-9}$ and $1.48 \times 10^{-8}$
    \item Exponential convergence with respect to collocation points and training parameters
    \item With 1\% noise, relative errors were around 0.026\% for $\alpha_1$ and 0.2\% for $\alpha_2$
\end{itemize}

\subsection{Helmholtz Equation with Inverse Variable Coefficient}

An inverse problem involving an unknown coefficient field:
\begin{align}
\frac{\partial^2 u}{\partial x^2} + \frac{\partial^2 u}{\partial y^2} - \gamma(x,y)u &= f(x,y) \\
\text{with boundary conditions and measurement data}
\end{align}

Key findings:
\begin{itemize}
    \item High accuracy with maximum $u$ error on the order of $10^{-8}$ and maximum $\gamma$ error on $10^{-5}$
    \item Exponential decrease in errors with respect to collocation points, training parameters, and measurement points
    \item For 1\% noise, the maximum relative error for $\gamma(x,y)$ was around 43\% (l$_\infty$ norm) and 5\% (l$_2$ norm)
    \item Regularization of output-layer coefficients improved accuracy in the presence of noise
\end{itemize}

\section{Implications and Advantages}

\subsection{Computational Performance}

\begin{itemize}
    \item Network training time grows approximately linearly with the number of collocation points and training parameters
    \item The method is significantly faster and more accurate than Physics-Informed Neural Networks (PINN) for inverse problems
    \item For noise-free data, the method achieves accuracy several orders of magnitude higher than PINN
\end{itemize}

\subsection{Robustness to Noise}

\begin{itemize}
    \item The method remains accurate even with noisy measurement data
    \item Scaling the measurement residual by a factor $\lambda_{mea}$ (0 < $\lambda_{mea}$ < 1) markedly improves accuracy for noisy data
    \item The robustness varies depending on the complexity of the inverse problem
\end{itemize}

\subsection{Practical Applications}

The method has significant implications for real-world applications where:
\begin{itemize}
    \item The physics is partially known (governing equations with unknown parameters)
    \item Only sparse measurements are available
    \item Parameter field distributions need to be identified
\end{itemize}

\subsection{Algorithm Selection}

The paper provides guidance on algorithm selection:
\begin{itemize}
    \item NLLSQ: Generally applicable to all problems, balance of efficiency and accuracy
    \item VarPro-F1: Effective for general cases, slightly more costly but can be more robust
    \item VarPro-F2: Most efficient for linear (with respect to solution) PDEs, but requires Newton iteration for nonlinear PDEs
\end{itemize}

The exponential convergence exhibited by the method suggests it could be particularly valuable for high-accuracy requirements in scientific computing and engineering applications.

\end{document} 